\section{Human validation of the intent criterion, in full}
\label{app:humanval}

The scale trend in \S\ref{sec:scale} rests on an automated classifier, so we validated it
against human judgement. 98 outputs were sampled \textbf{stratified by classifier verdict}
(40 \texttt{tight}, 28 \texttt{strong}-but-not-\texttt{tight}, 30 neither) --- deliberately
over-sampling the decision boundary so both false positives and false negatives are estimable.
Sampling is seeded and the annotator saw only the raw output, never the classifier's label.
We report \textbf{two rounds}, because the first is itself a finding.

\begin{center}
\small
\captionof{table}{Human validation of the intent criterion, round 1 and its adjudication. The classifier is used only to count what the server discarded, so its recall matters more than its precision.}\label{tab:validation}
\begin{tabular}{@{}lllllll@{}}
\toprule
& agreement & Cohen's $\kappa$ & precision & recall & FP & FN \\
\midrule
Round 1 (raw output only) & 86.7\% & 0.713 & 70.0\% & 96.6\% & 12 & 1 \\
After adjudication & \textbf{96.9\%} & \textbf{0.936} & \textbf{95.0\%} & 97.4\% & 2 & 1 \\
\bottomrule
\end{tabular}
\end{center}

The 12 round-1 false positives split into two causes, and the second is ours, not the
annotator's. \textbf{Eleven are reading-order failures}: the tool call sits mid-document, after
a fenced \verb|```python| block (match positions 1209--2585 in outputs of median length
${\sim}2000$); the annotator read the code block, concluded ``this is a direct answer, not a
call'', and stopped. \textbf{One is an instrument failure}: for a 32B item the matched span
begins at character 2585 and the annotation pack truncated every output at 2400, so the
evidence was \emph{not in the pack at all} and that item could not have been labelled
correctly. Re-shown the matched span alone, 10 of 13 disputed items were reversed. The
corrected pack carries every output in full.

\textbf{The first adjudication was one-sided, and we do not treat its $\kappa$ as a reliability
estimate.} Auditing the procedure after the fact: the 13 adjudicated items were selected
\emph{exactly} as the items where the annotator disagreed with the classifier (13/13 overlap),
and all 10 reversals moved \emph{toward} the classifier, none away. Under that selection rule
agreement can only rise, so $0.713 \to 0.936$ is a property of which items were re-examined,
not evidence of reliability. We report the second annotator instead.

\subsection*{A second independent annotator}

A second annotator, blind to the classifier, to the first annotator's labels and to the first
pack's item order, labelled the identical 98 items under a verbatim-identical rubric. Two
changes were made to the instrument: outputs are carried \textbf{in full}, and the adjudication
rule was fixed in advance to review \emph{all} three-way disagreements rather than only those
disagreeing with the classifier.

\begin{center}
\small
\captionof{table}{Inter-annotator reliability. Of 98 sampled items, 97 have complete paired A1/A2 records. The identical-byte subset removes a presentation confound introduced by the first pack's 2,400-character cap.}\label{tab:interannotator}
\begin{tabular}{@{}lcc@{}}
\toprule
Cohen's $\kappa$ (95\% CI, bootstrap) & all items ($n$=97) & \textbf{identical bytes ($n$=62)} \\
\midrule
A1 vs.\ A2 (inter-annotator) & 0.736 [0.590, 0.863] & \textbf{0.967 [0.897, 1.000]} \\
A2 vs.\ classifier           & 0.936 [0.855, 1.000] & \textbf{0.967 [0.897, 1.000]} \\
A1 vs.\ classifier           & 0.712 [0.559, 0.847] & 0.934 [0.834, 1.000] \\
\bottomrule
\end{tabular}
\end{center}

Of the 12 paired disagreements, 11 have full text exceeding 2,400 characters and A1 saw a
different presentation from A2; the first-round audit attributed most errors to reading order and
one to evidence actually omitted by truncation. We therefore do not assign all 11 to one cause.
On the 62 items where both read identical bytes the two agree on 61, and
\textbf{inter-annotator $\kappa$ = 0.967}.

\textbf{Third-party adjudication, and the final labels.} Fifteen items were not unanimous across
A1, A2 and the classifier. All fifteen went to a third adjudicator under the pre-registered rule
--- selected by \emph{any} three-way disagreement, never by disagreement with the classifier ---
who saw full text and no one else's labels. The adjudicator sided with A2 on 11, with the
classifier on 9, and with A1 on 6, and \textbf{moved away from the classifier on 6 of 15}. That
last number is the check that matters: the first round's adjudication moved away from the
classifier zero times out of thirteen, which is what an outcome looks like when the selection
rule has already decided it. Combining the 83 unanimous items with the 15 adjudicated ones gives
a 98-item gold standard. The third adjudicator was blind to classifier and annotator labels; the
classifier nevertheless affected the original stratified sample and which non-unanimous items
required adjudication, so we do not call the gold set classifier-independent.

\begin{center}
\small
\captionof{table}{Validation against the adjudicated gold standard. Pairwise denominators differ because one sampled item lacks a complete A1/A2 pair; classifier--gold uses all 98 items.}\label{tab:validation2}
\begin{tabular}{@{}lccc@{}}
\toprule
Cohen's $\kappa$ (95\% CI, bootstrap) & full $n$ & full sample & identical bytes ($n$=62) \\
\midrule
A1 vs.\ A2 (inter-annotator) & 97 & 0.736 [0.595, 0.866] & \textbf{0.967 [0.893, 1.000]} \\
Classifier vs.\ gold standard & 98 & \textbf{0.871 [0.756, 0.958]} & 0.934 [0.829, 1.000] \\
\bottomrule
\end{tabular}
\end{center}

\textbf{Correction factor.} Against the gold standard the \texttt{tight} stratum's precision is
36/40 = \textbf{0.900}, so \S\ref{sec:scale}'s headline 80/100 at 32B corrects to
\textbf{$\approx$72}. The criterion errs in both directions --- 4 over-counts against 2 calls it
misses entirely, a net over-count of 2 --- so the correction is smaller than the precision figure
alone suggests. \textbf{The trend is unaffected}: applying 0.900 uniformly gives
$0, 3.6, 18.9, 32.4, 72$ against a server-side zero at every size. We report uncorrected
classifier counts in all tables with this factor stated, rather than rescaling silently.

For the record, this number moved twice as the validation improved, and we report the trajectory
rather than only its endpoint: \textbf{0.950} under the first, one-sided adjudication;
\textbf{0.975} under A2's labels alone; \textbf{0.900} against the adjudicated gold standard. The
first was inflated by its selection rule, the second rested on a single rater, and only the third
is built from a rule fixed before the labels were seen.

\textbf{A bound on the measurement itself.} The probe stores at most 4000 characters of raw
output per item; 10 of the 98 sampled outputs reach that ceiling, and one grey-zone item is a
call truncated by it. The classifier and both annotators read identical bytes, which is what
agreement requires, but a call emitted beyond character 4000 is invisible to all three.
All ten cap hits occur in the 32B subset of this stratified sample: by increasing scale the
hit counts are 0/6, 0/13, 0/15, 0/24, and 10/40. These are sample diagnostics, not population
truncation rates, but they show that the uncertainty is concentrated at the headline rung.
The cap therefore creates downward pressure through false negatives. Because the classifier also
has false positives (tight precision 0.900), however, the net bias of the raw count is not known;
we report both the uncorrected count and the adjudication-based calibration rather than calling
the total a strict lower bound.
